\section{The Model}
\label{sec:model}

We consider a target thin region $A$ and an outside region $B$. A hub policy in region $A$ works on a \emph{local coordination} margin by raising the rate at which potentially valuable local interactions are realized. A pipeline policy works on a \emph{cross-regional renewal} margin by creating additional links between region $A$ and outside actors in region $B$. The model is dynamic only in a minimal two-period sense: it is designed to isolate one intertemporal welfare wedge, not to provide a full theory of cumulative regional path dependence.

The section also introduces the paper's principal distortion. The key wedge is a \emph{dynamic coordination externality}: decentralized actors value the immediate local matches generated by stronger coordination, but they do not fully internalize how dense current local interaction changes the common stock of future novelty available to the regional system. Repeated local interaction may exhaust novelty, while cross-regional renewal partially replenishes it.

\subsection{Environment and Policy Instruments}

Region $A$ contains $N_A\ge 2$ local participants, where a participant may be interpreted as a firm, researcher, entrepreneur, or other actor capable of generating collaborative surplus. Region $B$ contains $N_B\ge 1$ outside participants who form the external pool from which new ideas, collaborators, and partnerships may arrive. Region $B$ is not modeled as a symmetric second local economy. It is a reduced-form outside access margin through which cross-regional renewal can occur.

The number of unordered potential \emph{local} pairs within region $A$ is
\begin{equation}
P_A:=\frac{N_A(N_A-1)}{2},
\label{eq:pair_count}
\end{equation}
and the number of potential \emph{cross-regional} links between regions $A$ and $B$ is
\begin{equation}
Q_{AB}:=N_A N_B.
\label{eq:cross_pair_count}
\end{equation}

Two policy instruments are available. First, region $A$ may establish a public co-creation hub,
\[
h\in\{0,1\},
\]
where $h=1$ means that the hub is established. Second, the policymaker or decentralized local coalition may choose cross-regional pipeline effort
\[
b\ge 0.
\]
The variable $b$ summarizes brokerage, interregional outreach, visiting programs, and other efforts that connect region $A$ to outside actors.

Time is discrete and consists of two periods, $t=1,2$. The hub affects local matching in both periods. Pipeline effort affects only period 2.

\subsection{Local Coordination and Period-1 Surplus}

Let $m_A(h)$ denote the effective local matching rate among potential pairs within region $A$. We assume
\begin{equation}
m_A(h)=m_A^0+\bigl(m_A^1-m_A^0\bigr)h,
\qquad
0\le m_A^0<m_A^1\le 1.
\label{eq:matching_rate}
\end{equation}
Thus, establishing a hub lowers local search and meeting frictions in region $A$ and raises the fraction of potential local pairs that become realized economically meaningful interactions.

Let $v_L>0$ denote the average social surplus generated by one realized and still-novel local interaction. Period-1 gross surplus is therefore
\begin{equation}
B_1(h)=v_L P_A m_A(h).
\label{eq:B1}
\end{equation}

\subsection{Dynamic Coordination Externality}

The principal distortion in the model arises because the same local interactions that create value in period 1 can reduce the novelty available to the regional system in period 2. To capture this parsimoniously, let the period-2 \emph{local} component of social surplus be
\begin{equation}
B_{2L}(h,b)
=
\beta v_L P_A m_A(h)\bigl[1-\delta m_A(h)+\xi b\bigr],
\qquad
\beta\in(0,1],\;
\delta>0,\;
\xi>0.
\label{eq:B2_local}
\end{equation}

The term $\delta m_A(h)$ captures dynamic redundancy: more intensive local matching in period 1 raises the chance that period-2 interaction becomes repetitive and less novel. The term $\xi b$ captures the idea that cross-regional renewal partially replenishes the local opportunity set by bringing in outside ideas, collaborators, or perspectives.

We maintain the regularity condition
\begin{equation}
1-\delta m_A^1>0,
\label{eq:regularity}
\end{equation}
which ensures that second-period local novelty remains positive even without pipeline effort.

\subsection{Cross-Regional Links and a Minimally Active Outside Region}

Let the effective cross-regional connection intensity be
\begin{equation}
q(h,b)=q_0+\bigl(\eta_0+\eta_1 h\bigr)b,
\qquad
q_0\ge 0,
\quad
\eta_0>0,
\quad
\eta_1\ge 0.
\label{eq:q_b}
\end{equation}
The baseline term $q_0$ captures outside access that exists even without new policy. The parameter $\eta_0>0$ implies that pipeline effort can create value even without a hub: region $B$ is therefore not a purely passive outside pool. The complementarity term $\eta_1 h$ captures the idea that a hub in region $A$ can make brokerage more productive by organizing, absorbing, and channeling outside contacts.

Let $v_X>0$ denote the average surplus generated by one effective cross-regional linkage. The period-2 \emph{external} component of surplus is then
\begin{equation}
B_{2X}(h,b)
=
\beta v_X Q_{AB} q(h,b).
\label{eq:B2_external}
\end{equation}

Combining \eqref{eq:B2_local} and \eqref{eq:B2_external}, total period-2 gross surplus is
\begin{equation}
B_2(h,b)
=
\beta\Bigl\{
v_L P_A m_A(h)\bigl[1-\delta m_A(h)+\xi b\bigr]
+
v_X Q_{AB}\bigl[q_0+\bigl(\eta_0+\eta_1 h\bigr)b\bigr]
\Bigr\}.
\label{eq:B2}
\end{equation}

\subsection{Social Welfare, Contractibility, and the Decentralized Support Benchmark}

Establishing the hub requires a fixed cost $F>0$. Pipeline effort has convex cost
\[
\frac{\kappa}{2}b^2,
\qquad
\kappa>0.
\]

Total social welfare is therefore
\begin{equation}
W(h,b)
=
B_1(h)+B_2(h,b)-Fh-\frac{\kappa}{2}b^2.
\label{eq:W_general}
\end{equation}
Using \eqref{eq:B1} and \eqref{eq:B2}, this becomes
\begin{equation}
W(h,b)
=
v_L P_A m_A(h)
+
\beta\Bigl\{
v_L P_A m_A(h)\bigl[1-\delta m_A(h)+\xi b\bigr]
+
v_X Q_{AB}\bigl[q_0+\bigl(\eta_0+\eta_1 h\bigr)b\bigr]
\Bigr\}
-Fh-\frac{\kappa}{2}b^2.
\label{eq:W_expanded}
\end{equation}

To isolate the principal distortion, we introduce a decentralized support benchmark. Think of a local coalition, private operator, or decentralized set of sponsors that can contract on current hub services and on directly brokered outside links, but cannot fully contract on the continuation value of the regional opportunity set. The institutional logic is intentionally simple. Current meetings, memberships, and brokered introductions are observable and can therefore be priced. By contrast, once present interaction changes who knows whom inside the region, the resulting future novelty becomes a diffuse regional asset: later matches among third parties are hard to verify ex ante, exclusion is difficult, and much of the benefit takes the form of noncontractible option value. In that sense, the continuation value of future local novelty is treated as a local public-good or common-pool asset rather than as something owned by the current sponsoring coalition.

To keep the baseline sharp, the decentralized benchmark studies the polar case in which this continuation value is not internalized at all. Section~\ref{sec:extensions} relaxes that benchmark through partial internalization. Under this polar benchmark, the decentralized objective is
\begin{equation}
\Pi(h,b)
=
v_L P_A m_A(h)
+
\beta v_X Q_{AB}\bigl[q_0+\bigl(\eta_0+\eta_1 h\bigr)b\bigr]
-Fh-\frac{\kappa}{2}b^2.
\label{eq:Pi_general}
\end{equation}

The difference between \eqref{eq:W_expanded} and \eqref{eq:Pi_general} is exactly the common continuation value generated by the local regional system:
\[
\beta v_L P_A m_A(h)\bigl[1-\delta m_A(h)+\xi b\bigr].
\]
This is the paper's dynamic coordination externality in the baseline benchmark. Stronger local coordination creates private current benefits, but its effects on the future novelty of the local regional system are not fully priced by decentralized actors. The benchmark should therefore be read as a maintained appropriability structure, not as a claim that all decentralized settings literally involve zero internalization.

For later use, define
\begin{equation}
A_h^P
:=
v_X Q_{AB}\bigl(\eta_0+\eta_1 h\bigr),
\label{eq:Ah_private}
\end{equation}
and
\begin{equation}
A_h^W
:=
\xi v_L P_A m_A(h)+A_h^P.
\label{eq:Ah_social}
\end{equation}
Then \eqref{eq:Pi_general} and \eqref{eq:W_expanded} can be written as
\begin{equation}
\Pi(h,b)
=
\bar\Pi(h)+\beta A_h^P b-\frac{\kappa}{2}b^2,
\qquad
\bar\Pi(h):=v_L P_A m_A(h)+\beta v_X Q_{AB}q_0-Fh,
\label{eq:Pi_compact}
\end{equation}
and
\begin{equation}
W(h,b)
=
\bar W(h)+\beta A_h^W b-\frac{\kappa}{2}b^2,
\label{eq:W_compact}
\end{equation}
where
\begin{equation}
\bar W(h)
:=
v_L P_A m_A(h)
+
\beta\Bigl\{
v_L P_A m_A(h)\bigl[1-\delta m_A(h)\bigr]
+
v_X Q_{AB}q_0
\Bigr\}
-Fh.
\label{eq:Wbar}
\end{equation}

Since both \eqref{eq:Pi_compact} and \eqref{eq:W_compact} are strictly concave in $b$, the optimal pipeline choices conditional on $h$ are well defined. Section~\ref{sec:short_run} uses this benchmark to show why decentralized support and social evaluation need not coincide. Section~\ref{sec:hub_efficiency} first isolates the stand-alone hub comparison by setting $b=0$, and Section~\ref{sec:refresh} then restores pipeline effort to compare pipeline-alone, hub-alone, and bundled policy designs.
